\chapter{System Architecture}
\label{ch:arch}

Veilbox is designed as a layered system with clearly separated abstraction
boundaries. Each layer depends only on the layer below it and provides a
well-defined interface to the layer above. This chapter presents the
architecture at multiple levels of abstraction.

\section{Eight-Layer Architecture}

The system is organized into eight layers, from hardware at the bottom to
the user interface at the top. Each layer has a specific responsibility and
communicates with adjacent layers through standardized interfaces.

\subsection{Layer 1: Hardware}

The physical or virtual hardware provides the computing resources: CPU,
memory, storage, and networking. In a typical deployment, this is a QEMU/KVM
virtual machine with VirtIO paravirtualized devices. However, Veilbox boots
on any x86\_64 platform that provides:

\begin{itemize}[nosep]
  \item A BIOS or UEFI firmware interface.
  \item A bootable disk (IDE, SATA, VirtIO-BLK, NVMe).
  \item A network interface (VirtIO-NET, e1000, rtl8139).
  \item A serial port (16550A UART) for console access.
\end{itemize}

On bare metal, the same disk image boots via SeaBIOS or coreboot. The kernel
configured with VirtIO drivers will also work with AHCI/SATA and Intel
e1000/e1000e network controllers if those drivers are enabled.

\subsection{Layer 2: Firmware}

SeaBIOS (an open-source BIOS implementation shipped with QEMU) performs the
Power-On Self Test (POST) when the virtual machine starts:

\begin{enumerate}[nosep]
  \item CPU initialization: switches from real mode to protected mode briefly
        for memory sizing, then returns to real mode.
  \item Memory controller initialization: detects and sizes system RAM.
  \item PCI bus enumeration: discovers VirtIO devices, GPU (if present),
        and other PCI peripherals.
  \item Interrupt controller initialization: programs the PIC (8259A) and
        APIC (if available).
  \item Timer initialization: sets up the PIT (8254) and HPET.
  \item DMA controller initialization: programs the 8237 DMA controller.
  \item ACPI table generation: creates RSDT, DSDT, and other ACPI tables
        describing the virtual hardware configuration.
  \item Boot device scan: iterates through the boot order (set in QEMU's
        -boot option or the VM's BIOS menu), reading the first sector
        (LBA 0) of each device into memory at address 0x7C00.
\end{enumerate}

For each boot device, SeaBIOS reads the first 512 bytes into memory at
address 0x7C00 using INT 13h function 42h (extended read). It then checks
whether the last two bytes at 0x7DFE--0x7DFF are 0x55 0xAA (the MBR boot
signature). If the signature is present, SeaBIOS sets the DL register to
the boot drive number and transfers control by executing a far jump to
0x7C00:0000.

\subsection{Layer 3: Bootloader}

GRUB (Grand Unified Bootloader) version 2.x runs in BIOS mode (i386-pc
platform). It consists of three stages:

\paragraph{Stage 1: boot.img}
This 512-byte file is written to LBA 0 (the Master Boot Record) by the
Veilbox GRUB installer. It contains:

\begin{itemize}[nosep]
  \item A short real-mode assembly program (~446 bytes) that sets up segment
        registers, initializes the stack, and reads the blocklist from a
        fixed offset in the file.
  \item The partition table (64 bytes at offset 0x1BE--0x1FD) describing
        the disk layout.
  \item The boot signature (2 bytes at offset 0x1FE--0x1FF = 55 AA).
\end{itemize}

When SeaBIOS jumps to 0x7C00, boot.img executes and reads the blocklist
(embedded in the core.img file at offset 0x1F4) to determine how many
sectors to load and where to place them in memory.

\paragraph{Stage 1.5: core.img}
This file is written to the sectors immediately following the MBR (LBA 1
onwards). It contains:

\begin{itemize}[nosep]
  \item diskboot.img: the first sector of core.img, which loads the rest of
        core.img. It switches from real mode to protected mode, enables the
        A20 gate (allowing access to memory above 1 MB), and sets up the
        protected-mode environment for the GRUB kernel.
  \item Filesystem drivers: ext2 (for reading the VEILBOX partition),
        part\_msdos (for parsing the MBR partition table), and biosdisk
        (for BIOS INT 13h disk access).
  \item The GRUB kernel: a small embedded operating system that provides
        command parsing, device management, and module loading.
  \item An embedded configuration: added by the -c option to
        grub2-mkimage, this tells GRUB to search for the VEILBOX partition
        by filesystem label and set the prefix path.
\end{itemize}

\paragraph{Stage 2: /boot/grub/}
GRUB modules and configuration files on the VEILBOX ext4 partition. The
embedded config in core.img executes "search --label --set=root VEILBOX"
to find the partition, then "set prefix=(\$root)/boot/grub" to locate the
GRUB directory. GRUB reads grub.cfg, displays the boot menu, and loads
the kernel.

\subsection{Layer 4: Kernel}

Linux kernel 7.1.0-rc6 compiled with custom configuration for virtualized
environments. The kernel binary (vmlinuz) includes the embedded initramfs
and is approximately 95 MB in size.

Key configuration areas:
\begin{itemize}[nosep]
  \item VirtIO drivers for paravirtualized block and network I/O.
  \item Full namespace and cgroup support for container isolation.
  \item Seccomp and BPF for syscall filtering.
  \item Serial console support for headless operation.
  \item EFI stub for QEMU -kernel compatibility.
  \item Striped-down driver set (no sound, USB, Wi-Fi, etc.).
\end{itemize}

\subsection{Layer 5: Init System}

BusyBox init runs as PID 1, the first userspace process created by the
kernel. It reads /etc/inittab and executes the configured actions:

\begin{enumerate}[nosep]
  \item sysinit: run /etc/init.d/rcS to mount filesystems and start services.
  \item respawn: start getty (or autologin) on tty1 and ttyS0, restarting
        them whenever they exit.
  \item shutdown: unmount filesystems and deactivate swap when the system
        is told to power off or reboot.
\end{enumerate}

\subsection{Layer 6: Boot Scripts}

The /etc/init.d/rcS script is the single point of customization for system
behavior. It performs:

\begin{enumerate}[nosep]
  \item Mount virtual filesystems: proc, sysfs, devtmpfs, devpts, tmpfs.
  \item Set hostname from /etc/hostname.
  \item Mount the persistent state partition with a multi-tier fallback
        strategy: external VirtIO disk (vdb), external SATA disk (sdb),
        VEILBOX partition /state directory, or memory-backed tmpfs.
  \item Bring up network interfaces: loopback (lo) and eth0.
  \item Start the DHCP client (udhcpc) in background.
  \item Start the system logger (syslogd) in background.
  \item Start containerd in background.
  \item Start Dropbear SSH server in foreground.
\end{enumerate}

\subsection{Layer 7: Services}

Four services run in the background after rcS completes:

\begin{itemize}[nosep]
  \item \textbf{udhcpc}: BusyBox DHCP client. Negotiates an IP address from
        the QEMU SLiRP virtual DHCP server. The script at
        /usr/share/udhcpc/default.script configures the interface on
        bound/renew events.
  \item \textbf{syslogd}: BusyBox system logger. Captures kernel log messages
        (via /proc/kmsg) and process messages (via /dev/log) to
        /var/log/messages.
  \item \textbf{containerd}: Container runtime daemon. Exposes a gRPC API at
        /run/containerd/containerd.sock for container lifecycle management.
  \item \textbf{Dropbear}: Lightweight SSH server. Listens on port 22 for
        incoming connections. Supports key-based and password authentication.
\end{itemize}

\subsection{Layer 8: User Interface}

The login prompt is presented on two consoles:
\begin{itemize}[nosep]
  \item \textbf{tty1}: VGA text console (available when QEMU is run with a
        graphical window or when booting on bare metal with a monitor).
  \item \textbf{ttyS0}: Serial console (available via QEMU's -nographic
        mode or a physical serial port).
\end{itemize}

The autologin wrapper (/sbin/autologin) checks /proc/cmdline for the
veilbox.autologin parameter. If present, it runs "login -f root" to
automatically log in as root. Otherwise, it runs getty to present the
standard login prompt.

\section{The Build Pipeline}

The operating system is assembled by build.sh in a linear pipeline of
14 stages. Each stage is independent and can be run separately for
debugging or incremental development.

\begin{enumerate}[nosep]
  \item \textbf{Install dependencies}: Check for required packages, install
        via sudo dnf if on Fedora.
  \item \textbf{Create rootfs directories}: Set up the Linux filesystem
        hierarchy (bin, sbin, etc, usr, lib, etc.).
  \item \textbf{Write configuration files}: Create inittab, rcS, passwd,
        shadow, containerd config, DHCP script, profile.
  \item \textbf{Configure kernel}: Run make defconfig, merge custom fragment,
        resolve dependencies with olddefconfig.
  \item \textbf{Strip unnecessary drivers}: Disable sound, USB, Wi-Fi,
        Bluetooth, media, ATA, SCSI, and other VM-unneeded subsystems.
  \item \textbf{Populate rootfs binaries}: Download and install BusyBox,
        containerd, runc, nerdctl, and Dropbear.
  \item \textbf{Generate SSH keys}: Create Ed25519 key pair for root access.
  \item \textbf{Copy shared libraries}: Find dynamically linked binaries and
        copy their library dependencies to /lib64.
  \item \textbf{Generate initramfs}: Create the kernel file list from the
        rootfs directory, add device nodes, and set as initramfs source.
  \item \textbf{Build kernel}: Compile the kernel with make -j\$(nproc)
        and copy the resulting bzImage to the output directory.
  \item \textbf{Create SquashFS reference}: Compress the rootfs for
        inspection (optional, not used at boot time).
  \item \textbf{Create state image}: Make a 128 MB ext4 filesystem for
        legacy persistent storage.
  \item \textbf{Create bootable disk}: Build the 8 GB MBR-partitioned disk
        image with ext4 filesystem, kernel, GRUB config, and rootless
        GRUB installation.
  \item \textbf{Convert to VDI}: Use qemu-img to create a VirtualBox-
        compatible sparse disk image.
\end{enumerate}

\section{Design Trade-offs}

\subsection{Embedded Initramfs vs. Separate Initrd}

The decision to embed the initramfs inside the kernel binary (rather than
keeping it as a separate file on the boot partition) has far-reaching
consequences for the system architecture.

\textbf{Advantages:}
\begin{itemize}[nosep]
  \item Single-file deployment: copy one file (vmlinuz) to the boot
        partition and the system is ready.
  \item Atomic kernel+rootfs: the kernel and its root filesystem are always
        in sync. There is no risk of a kernel update that is incompatible
        with an older initramfs.
  \item Bootloader simplicity: the bootloader loads one file (vmlinuz)
        instead of two (kernel + initrd). The GRUB configuration is simply
        "linux /boot/vmlinuz console=ttyS0".
  \item QEMU compatibility: the -kernel option loads the kernel and
        automatically detects the embedded initramfs, eliminating the
        need for -initrd.
\end{itemize}

\textbf{Disadvantages:}
\begin{itemize}[nosep]
  \item Larger kernel binary: ~95 MB vs. ~10 MB for a bare kernel. This
        increases boot time (decompression of the initramfs) and storage
        requirements.
  \item Longer build time: every change to the rootfs requires a kernel
        rebuild. The make process re-embeds the initramfs each time.
  \item No runtime rootfs changes: the root filesystem cannot be modified
        without rebuilding the kernel.
\end{itemize}

\subsection{BusyBox vs. Full Userspace}

BusyBox provides over 300 Unix utilities in a single 1 MB binary. A full
GNU userspace with coreutils, util-linux, procps-ng, bash, and other
standard packages would be 200+ MB.

\textbf{Advantages:}
\begin{itemize}[nosep]
  \item Dramatically smaller disk and memory footprint.
  \item Simpler dependency management (one binary, one set of libraries).
  \item Faster boot (fewer binaries to load).
  \item Smaller attack surface (fewer lines of code).
\end{itemize}

\textbf{Disadvantages:}
\begin{itemize}[nosep]
  \item Limited tool selection: many GNU extensions are not available.
  \item Reduced functionality: ps has fewer options, grep is slower on
        large files, awk is not included (use the smaller BusyBox awk).
  \item Different error messages and behavior: scripts written for GNU
        tools may need adjustment.
\end{itemize}

\subsection{BIOS vs. UEFI Boot}

Veilbox uses GRUB in BIOS/CSM mode (i386-pc) rather than UEFI (x86\_64-efi).

\textbf{Advantages:}
\begin{itemize}[nosep]
  \item Maximum compatibility with older firmware and hypervisors.
  \item Simpler installer (boot.img + core.img at known LBA positions).
  \item Works with QEMU's default SeaBIOS without UEFI firmware.
\end{itemize}

\textbf{Disadvantages:}
\begin{itemize}[nosep]
  \item No Secure Boot support.
  \item Limited to 2 TB boot disks (MBR limitation).
  \item Slower boot (real-mode transitions).
\end{itemize}
